% Copyright (c) 2026 Dennis Michael Heine.
% Released under the CC BY-NC-SA 4.0 license as described in the file LICENSE.
\documentclass[11pt]{amsart}
\usepackage[T1]{fontenc}
\usepackage{amsmath,amssymb,amsthm}
\usepackage{listings}
\usepackage[colorlinks=true,linkcolor=blue,citecolor=blue,urlcolor=blue]{hyperref}

\lstset{basicstyle=\ttfamily\small,breaklines=true,columns=fullflexible,
  keepspaces=true,frame=single,framesep=4pt}

\newtheorem{theorem}{Theorem}[section]
\newtheorem{lemma}[theorem]{Lemma}
\newtheorem{definition}[theorem]{Definition}
\newtheorem{remark}[theorem]{Remark}

\title[A machine-checked Dobrushin criterion]{A machine-checked Dobrushin
criterion for abstract polymer systems}
\author{Dennis Michael Heine}
\date{September 3, 2026. \emph{Preliminary draft.}}

\begin{document}

\begin{abstract}
We describe a formalisation, in Lean~4 over \textsf{mathlib}, of the core of
the cluster expansion method for abstract polymer systems. Two results are
proved in full, with no unproved placeholders: the deletion recursion
$Z(\Lambda) = Z(\Lambda \setminus \{\gamma\}) + w(\gamma)\,
Z(\Lambda \setminus N^*(\gamma))$ for the polymer partition function, for
weights in an arbitrary commutative ring; and, for real weights, the
Dobrushin convergence criterion in product form: if
$\lvert w(\gamma)\rvert \prod_{\delta \nsim \gamma} (1 + \mu(\delta)) \le
\mu(\gamma)$ for some $\mu \ge 0$, then $Z(\Lambda) \neq 0$ and
$\lvert Z(\Lambda \setminus \{\gamma\})\rvert \le (1 + \mu(\gamma))
\lvert Z(\Lambda)\rvert$. We record a point that the formalisation forced
into the open: the ratio-telescoping induction that proves the product-form
criterion does not prove the classical Kotecký--Preiss criterion in sum
form, whose formalisation via the cluster-tree induction remains future
work. The development is a first step of a machine-checked treatment of the
convergence machinery underlying rigorous renormalisation group analyses of
lattice field theories.
\end{abstract}

\maketitle

\section{Introduction}

Cluster expansions are among the basic tools of rigorous statistical
mechanics and constructive field theory. In the abstract polymer formulation
one is given a finite set of \emph{polymers}, a symmetric and reflexive
\emph{incompatibility} relation, and complex or real \emph{weights}; the
object of interest is the polymer partition function, a sum over collections
of pairwise compatible polymers. The Kotecký--Preiss criterion
\cite{KP86} gives a sufficient condition for the logarithm of this partition
function to be controlled by an absolutely convergent series, uniformly in
the volume; see \cite{FV17} for a textbook treatment and \cite{FP07} for a
sharper criterion.

Two circumstances motivate a formalisation. First, applications of the
method, notably in rigorous renormalisation group analyses of lattice gauge
theories in the tradition of Balaban (see \cite{Dim13a,Dim13b,Dim13c} for a
modern exposition), rest on long chains of estimates whose unnamed constants
have, as far as we are aware, never been verified line by line by a machine;
the Liquid Tensor Experiment \cite{LTE} has demonstrated that formalisation
of proofs at this level of technical difficulty is feasible and useful.
Second, the library situation is favourable: \textsf{mathlib}
\cite{mathlib20} contains the finite set combinatorics, ordered field and
special function material needed for the statement of the criterion, and
recent formalisations in stochastic analysis \cite{BM25} and high energy
physics tooling \cite{HepLean} indicate growing infrastructure on the
mathematical physics side. To our knowledge no cluster expansion material is
currently available in \textsf{mathlib}.

The present note documents the first two steps. First, a machine-checked
proof of the fundamental recursion for the polymer partition function, in a
generality (commutative ring valued weights) that covers the real, complex
and formal power series cases occurring in applications. Second, a complete
proof of a convergence criterion: the Dobrushin product-form condition
implies nonvanishing of $Z$ together with the ratio bound that controls the
effect of deleting a polymer. The development is about 400 lines of Lean,
compiles against current \textsf{mathlib}, and contains no \texttt{sorry}
placeholders. The classical Kotecký--Preiss criterion in sum form is stated
as a definition; the formalisation made unusually explicit why its proof
requires a different induction, as discussed in Section~\ref{sec:kp}.

\section{The abstract setting}

\begin{definition}
A \emph{polymer system} on a type $\iota$ consists of a relation
$\iota \times \iota \to \{\mathrm{true},\mathrm{false}\}$, written
$\gamma \nsim \delta$, that is symmetric and satisfies
$\gamma \nsim \gamma$ for every $\gamma$. A finite set
$S \subseteq \iota$ is \emph{independent} if $\gamma \sim \delta$ for all
distinct $\gamma, \delta \in S$.
\end{definition}

\begin{definition}
Let $R$ be a commutative ring, $w \colon \iota \to R$, and let
$\Lambda \subseteq \iota$ be finite. The \emph{partition function} is
\[
  Z(\Lambda) \;=\; \sum_{\substack{S \subseteq \Lambda \\ S \text{
  independent}}} \; \prod_{\gamma \in S} w(\gamma),
\]
the empty set contributing the term $1$.
\end{definition}

For $\gamma \in \Lambda$ write
$N^*(\gamma) = \{\delta \in \Lambda : \delta \nsim \gamma\}$ for the closed
incompatibility neighbourhood, so that $\gamma \in N^*(\gamma)$, and
$\Lambda_\gamma = \Lambda \setminus N^*(\gamma)$ for the set of polymers in
$\Lambda$ compatible with $\gamma$.

\begin{theorem}[fundamental recursion]\label{thm:rec}
For every finite $\Lambda$ and every $\gamma \in \Lambda$,
\[
  Z(\Lambda) \;=\; Z(\Lambda \setminus \{\gamma\})
  \;+\; w(\gamma)\, Z(\Lambda_\gamma).
\]
\end{theorem}

\begin{proof}
Split the sum defining $Z(\Lambda)$ according to whether $\gamma \in S$.
The independent sets not containing $\gamma$ are exactly the independent
subsets of $\Lambda \setminus \{\gamma\}$. The independent sets containing
$\gamma$ are of the form $S = \{\gamma\} \cup T$ with
$T \subseteq \Lambda_\gamma$ independent, and this correspondence
$S \mapsto S \setminus \{\gamma\}$ is a bijection; the associated products
factor as $w(\gamma) \prod_{\delta \in T} w(\delta)$.
\end{proof}

The recursion is the engine of the ratio-telescoping induction, which
yields the following convergence criterion, stated here in the product
(Dobrushin) form; see \cite{SS05} for the surrounding theory.

\begin{theorem}[Dobrushin criterion; formalised]\label{thm:dob}
Let $w, \mu \colon \iota \to \mathbb{R}$ with $\mu \ge 0$ on $\Lambda$, and
suppose that for every $\gamma \in \Lambda$
\[
  \lvert w(\gamma)\rvert \prod_{\substack{\delta \in \Lambda \\ \delta
  \nsim \gamma}} \bigl(1 + \mu(\delta)\bigr) \;\le\; \mu(\gamma),
\]
the product running over the closed incompatibility neighbourhood. Then
$Z(\Lambda) \neq 0$, and for every $\gamma \in \Lambda$,
\[
  \lvert Z(\Lambda \setminus \{\gamma\})\rvert \;\le\;
  \bigl(1 + \mu(\gamma)\bigr)\, \lvert Z(\Lambda)\rvert .
\]
\end{theorem}

\begin{proof}[Proof (as formalised)]
Strong induction on $\lvert\Lambda\rvert$; the two statements are proved
simultaneously and no division occurs anywhere. Given $\gamma \in \Lambda$,
one first shows by an inner induction that deleting any set $E$ from a
proper subset $S$ costs at most the factor $\prod_{\delta \in E}(1 +
\mu(\delta))$, each single deletion being controlled by the outer induction
hypothesis. Applying this with $S = \Lambda \setminus \{\gamma\}$ and $E$
the incompatibility neighbourhood of $\gamma$ bounds
$\lvert Z(\Lambda_\gamma)\rvert$; the recursion of Theorem~\ref{thm:rec}
and the triangle inequality then give
$\lvert Z(\Lambda \setminus \{\gamma\})\rvert \le \lvert Z(\Lambda)\rvert +
c\,\lvert Z(\Lambda \setminus \{\gamma\})\rvert$ with
$c = \lvert w(\gamma)\rvert \prod_{\delta \in D}(1 + \mu(\delta))$, and the
hypothesis, rewritten as $c\,(1 + \mu(\gamma)) \le \mu(\gamma)$ after
splitting off the factor for $\delta = \gamma$, closes the estimate.
Nonvanishing follows by applying the ratio bound at any single polymer.
\end{proof}

\section{The formalisation}

The development lives in a single file \texttt{ClusterExpansion.lean} inside
a Lean~4 project depending on \textsf{mathlib}. We record the design
decisions that shaped it.

\subsection{Decidability by construction}
The incompatibility relation is carried as a \texttt{Bool}-valued function
rather than a \texttt{Prop}-valued one:

\begin{lstlisting}
structure PolymerSystem (i : Type*) where
  incomp : i -> i -> Bool
  symm   : forall g d, incomp g d = incomp d g
  refl   : forall g, incomp g g = true
\end{lstlisting}

Independence is then a decidable predicate for free, which is what permits
the partition function to be written as a sum over a filtered powerset,
\begin{lstlisting}
def Z (L : Finset i) : R :=
  Sum_{S in L.powerset.filter (fun S => Indep P S)}
    Prod_{g in S} w g
\end{lstlisting}
(displayed here in mathematical notation; the source uses the
\textsf{mathlib} big operator syntax). No classical choice is needed to
\emph{state} anything; \texttt{classical} is invoked only inside the main
proof for convenience.

\subsection{Generality of the weight ring}
Theorem~\ref{thm:rec} is purely combinatorial, and the formal statement
quantifies over an arbitrary \texttt{CommRing}. This matters for the
intended applications: cluster expansions are routinely applied with complex
weights and with weights in rings of formal power series, and the analytic
hypotheses enter only later, through the convergence criteria.

\subsection{The bijection}
The formal proof follows the paper proof. The sum is split with
\texttt{sum\_filter\_add\_sum\_filter\_not}; the first half is identified
with $Z(\Lambda \setminus \{\gamma\})$ by an extensionality argument on the
filtered powersets, using
$S \subseteq \Lambda \setminus \{\gamma\} \leftrightarrow
S \subseteq \Lambda \wedge \gamma \notin S$. The second half is treated with
\texttt{Finset.sum\_bij'}, the bijection being
$S \mapsto S \setminus \{\gamma\}$ with inverse
$T \mapsto \{\gamma\} \cup T$. The single nontrivial lemma feeding the
bijection is:

\begin{lstlisting}
theorem indep_insert_iff (hT : T <= compat P L g) :
    Indep P (insert g T) <-> Indep P T
\end{lstlisting}

\noindent
whose proof is a case analysis using symmetry of the relation. The product
identity $\prod_{S} w = w(\gamma)\prod_{S \setminus \{\gamma\}} w$ is
\texttt{Finset.prod\_insert} applied after rewriting $S$ as
\texttt{insert} $\gamma$ \texttt{(S.erase} $\gamma$\texttt{)}.

\section{Why the sum form is not a corollary}\label{sec:kp}

The classical Kotecký--Preiss criterion \cite{KP86} is stated in sum form:
$a \ge 0$ and, for every $\gamma \in \Lambda$,
\[
  \sum_{\substack{\delta \in \Lambda \\ \delta \nsim \gamma}}
  \lvert w(\delta)\rvert \, e^{a(\delta)} \;\le\; a(\gamma).
\]
It is tempting to expect Theorem~\ref{thm:dob} to yield this by the
substitution $\mu = e^{a} - 1$; the formalisation forced us to notice that
it does not. The telescoping argument produces \emph{products} of the
factors $1 + \mu(\delta)$ over a neighbourhood, and the sum-form hypothesis
bounds a \emph{sum} over the same neighbourhood; there is no pointwise
implication in either direction. A two-polymer example makes this concrete:
for a pair of mutually incompatible polymers with equal weights $w$ and
equal $a$, the sum condition permits $w = a e^{-a}/2$, while the product
condition at $\mu = e^{a} - 1$ demands $a/2 \le 1 - e^{-a}$, which fails for
$a$ beyond roughly $1.59$. The two criteria are genuinely different
theorems. The sum form is proved in the literature by an induction over
clusters with tree-graph bounds (see \cite[Chapter~5]{FV17}); formalising
that induction, and then the convergent expansion for $\log Z$ itself, with
the Fernández--Procacci refinement \cite{FP07} as a further target, is the
next stage of this programme. In the development the sum-form condition is
present as a definition, deliberately without an accompanying unproved
theorem.

\section{Related work}

We are not aware of prior formalisations of cluster expansion material in
any proof assistant. On the neighbouring shelves, \textsf{mathlib}
\cite{mathlib20} supplies the combinatorial and analytic prerequisites; the
recent formalisation of Brownian motion \cite{BM25} built Gaussian measure
and Kolmogorov extension infrastructure relevant to constructive field
theory; and the PhysLean effort \cite{HepLean} is assembling high energy
physics definitions in Lean. The exposition of Balaban's renormalisation
group method by Dimock \cite{Dim13a,Dim13b,Dim13c} is the natural blueprint
document for the larger programme, in the sense of the blueprint-driven
formalisations pioneered around \cite{LTE}.

\section{Availability}

The development (project \texttt{KPLean}, file
\texttt{KPLean/ClusterExpansion.lean}, Lean toolchain \texttt{v4.33.1})
builds with \texttt{lake build} against the \textsf{mathlib} version pinned
in the accompanying manifest. It is currently distributed as part of a
larger repository of the author; extraction into a standalone repository is
planned.

\subsection*{Acknowledgements}
The formalisation and this note were produced with substantial assistance
from an AI system (Claude, Anthropic), including proof drafting and
literature checks; all remaining errors are the author's.

\begin{thebibliography}{9}

\bibitem{KP86} R.~Kotecký and D.~Preiss, \emph{Cluster expansion for
abstract polymer models}, Comm. Math. Phys. \textbf{103} (1986), 491--498.

\bibitem{FV17} S.~Friedli and Y.~Velenik, \emph{Statistical Mechanics of
Lattice Systems: A Concrete Mathematical Introduction}, Cambridge University
Press, 2017.

\bibitem{FP07} R.~Fernández and A.~Procacci, \emph{Cluster expansion for
abstract polymer models: new bounds from an old approach}, Comm. Math. Phys.
\textbf{274} (2007), 123--140.

\bibitem{SS05} A.~D.~Scott and A.~D.~Sokal, \emph{The repulsive lattice
gas, the independent-set polynomial, and the Lovász local lemma}, J. Stat.
Phys. \textbf{118} (2005), 1151--1261.

\bibitem{Dim13a} J.~Dimock, \emph{The renormalization group according to
Balaban I: small fields}, Rev. Math. Phys. \textbf{25} (2013), 1330010;
arXiv:1108.1335.

\bibitem{Dim13b} J.~Dimock, \emph{The renormalization group according to
Balaban II: large fields}, J. Math. Phys. \textbf{54} (2013), 092301;
arXiv:1212.5562.

\bibitem{Dim13c} J.~Dimock, \emph{The renormalization group according to
Balaban III: convergence}, Ann. Henri Poincaré \textbf{15} (2014),
2133--2175; arXiv:1304.0705.

\bibitem{LTE} J.~Commelin and A.~Topaz, \emph{Abstraction boundaries and
spec driven development in pure mathematics}, Bull. Amer. Math. Soc.
\textbf{61} (2024), 241--255; see also the Liquid Tensor Experiment,
\url{https://leanprover-community.github.io/blog/posts/lte-final/}.

\bibitem{mathlib20} The mathlib Community, \emph{The Lean mathematical
library}, in: Proc. 9th ACM SIGPLAN Int. Conf. on Certified Programs and
Proofs (CPP 2020), 367--381.

\bibitem{BM25} R.~Degenne, P.~Ledvinka, E.~Marion and P.~Pfaffelhuber,
\emph{A formalization of Brownian motion in Lean}, arXiv:2511.20118.

\bibitem{HepLean} J.~Tooby-Smith, \emph{HepLean: digitalising high energy
physics}, arXiv:2405.08863.

\end{thebibliography}

\end{document}
